\clearpage
\VerdierPhysicalPageAnchor{0064}%
\VerdierSemanticAnchor{ch1-1-6-6-functor-characterization}
由命题 (1.5.2)，给定一个从 $C$ 到 $D$ 的函子，等价于给定一个函子

\VerdierSemanticAnchor{ch1-1-6-6-eq2}
\begin{equation}
  (X_1,\ldots,X_n)\longmapsto F(X_1,\ldots,X_n)
  \tag{1.6.6.2}
\end{equation}

\VerdierSemanticAnchor{ch1-1-6-6-transition-isomorphisms}
以及函子同构

\VerdierSemanticAnchor{ch1-1-6-6-eq3}
\begin{equation}
  p_i\colon
  F(X_1,\ldots,X_i[1],\ldots,X_n)
  \xrightarrow{\ \sim\ }
  F(X_1,\ldots,X_n)[1]
  \tag{1.6.6.3}
\end{equation}

\VerdierSemanticAnchor{ch1-1-6-6-anticommutativity}
并要求图 (1.5.1.2)

\VerdierSemanticAnchor{ch1-1-6-6-anticommutativity-diagram}
\begin{equation*}
  \resizebox{\textwidth}{!}{%
    $\begin{tikzcd}[ampersand replacement=\&,column sep=2.8em,row sep=3.8em]
      F(X_1,\ldots,X_i[1],\ldots,X_j[1],\ldots,X_n)
        \arrow[r,"{p_i}"]
        \arrow[d,"{p_j}"]
        \& F(X_1,\ldots,X_i,\ldots,X_j[1],\ldots,X_n)[1]
          \arrow[d,"{p_j[1]}"] \\
      F(X_1,\ldots,X_i[1],\ldots,X_j,\ldots,X_n)[1]
        \arrow[r,"{p_i[1]}"]
        \& F(X_1,\ldots,X_i,\ldots,X_j,\ldots,X_n)[2]
    \end{tikzcd}$%
  }
\end{equation*}

对 $i<j$ 是反交换的。（为简化图 (1.5.1.2)，我们假定 $D$ 的过渡同构为
恒等态射；本书中始终如此。）

\VerdierSemanticAnchor{ch1-1-6-6-morphism-characterization}
同样，由命题 (1.5.10)，我们所考察的函子态射
$u\colon F\to F'$ 必须满足以下条件：下图交换：

\VerdierSemanticAnchor{ch1-1-6-6-eq4}
\begin{equation}
  \resizebox{0.92\textwidth}{!}{%
    $\begin{tikzcd}[ampersand replacement=\&,column sep=5.0em,row sep=3.8em]
      F(X_1,\ldots,X_i[1],\ldots,X_n)
        \arrow[r,"u"]
        \arrow[d,"{p_i}"]
        \& F'(X_1,\ldots,X_i[1],\ldots,X_n)
          \arrow[d,"{p'_i}"] \\
      F(X_1,\ldots,X_i,\ldots,X_n)[1]
        \arrow[r,"{u[1]}"]
        \& F'(X_1,\ldots,X_i,\ldots,X_n)[1]
    \end{tikzcd}$%
  }
  \tag{1.6.6.4}
\end{equation}

\VerdierSemanticAnchor{ch1-1-6-6-underlying-convention}
当上下文不会造成混淆时，我们用底层范畴表示 $\mathbb Z$-范畴，并用底层函子
表示扭曲的 $\mathbb Z^n$-函子。
